\section{Acyclic resolutions compute right-derived functors}
\label{mk-chap:Cohomology_AcyclicResolution}

\section{Known results}

The following three well-known results underpin the injective-resolution
construction of \(R^n G\) and are the foundational inputs to the argument.

\begin{lemma}[Right-derived functor via an injective resolution]
  \label{mk-lem:right_derived_injective_resolution}
  \lean{CategoryTheory.InjectiveResolution.isoRightDerivedObj}
  \mathlibok
  \dcref{custom}

  Let \(\mathcal{A}\) be an abelian category with enough injectives,
  \(\mathcal{B}\) abelian, and \(G : \mathcal{A} \to \mathcal{B}\) additive. For an
  object \(X\) of \(\mathcal{A}\), a choice of injective resolution
  \(I\) of \(X\) (an exact complex \(0 \to X \to I^0 \to I^1 \to \cdots\) with each
  \(I^n\) injective), and every \(n \in \mathbb{N}\), there is a canonical
  isomorphism in \(\mathcal{B}\)
  \[
    (R^n G)(X)
      \;\cong\;
      H^n\bigl(G(I^\bullet)\bigr),
  \]
  where \(G(I^\bullet)\) is the cochain complex obtained by applying \(G\)
  degreewise to the resolution complex \(I^\bullet\) and \(H^n\) is its \(n\)-th
  cohomology object. This is the defining computation of the right-derived functor
  and the seed from which the acyclic-resolution generalization is bootstrapped.
\end{lemma}

\begin{lemma}[Higher derived functors vanish on injectives]
  \label{mk-lem:right_derived_vanishes_injective}
  \lean{CategoryTheory.Functor.isZero_rightDerived_obj_injective_succ}
  \mathlibok
  \dcref{custom}

  With \(\mathcal{A}, \mathcal{B}, G\) as above, if \(J\) is an injective object of
  \(\mathcal{A}\) then \((R^{n+1} G)(J) = 0\) for every \(n \in \mathbb{N}\); that
  is, the higher right-derived functors vanish on injective objects. In particular
  every injective object is right-\(G\)-acyclic in the sense of
  Definition~\ref{mk-def:right_acyclic}.
\end{lemma}

\begin{lemma}[The zeroth derived functor of a left-exact functor]
  \label{mk-lem:right_derived_zero_iso_self}
  \lean{CategoryTheory.Functor.rightDerivedZeroIsoSelf}
  \mathlibok
  \dcref{custom}

  With \(\mathcal{A}, \mathcal{B}, G\) as above, if \(G\) is left exact (preserves
  finite limits) then the natural transformation \(G \to R^0 G\) is an isomorphism,
  \[
    R^0 G \;\cong\; G.
  \]
\end{lemma}

\begin{lemma}[Long exact homology sequence of a short exact sequence of complexes]
  \label{mk-lem:homology_long_exact_sequence}
  \lean{CategoryTheory.ShortComplex.ShortExact.homology_exact₁,
    CategoryTheory.ShortComplex.ShortExact.homology_exact₂,
    CategoryTheory.ShortComplex.ShortExact.homology_exact₃,
    CategoryTheory.ShortComplex.ShortExact.δ}
  \mathlibok
  \dcref{custom}

  Let \(\mathcal{C}\) be an abelian category and let
  \[
    0 \longrightarrow S.X_1 \longrightarrow S.X_2 \longrightarrow S.X_3
      \longrightarrow 0
  \]
  be a short exact sequence of cochain complexes in \(\mathcal{C}\) --- that is, a
  short complex \(S = (S.X_1 \to S.X_2 \to S.X_3)\) of objects of the category of
  \(\mathbb{Z}\)-indexed cochain complexes \(\mathrm{HomologicalComplex}\,\mathcal{C}\)
  which is short exact (degreewise mono, epi and exact in the middle). Then for every
  degree \(n\) there is a connecting morphism
  \[
    \delta^n \;:\; H^n(S.X_3) \longrightarrow H^{n+1}(S.X_1)
  \]
  (the Mathlib map \texttt{ShortComplex.ShortExact.\(\delta\)}), and these assemble
  into a long exact sequence of homology objects in \(\mathcal{C}\)
  \[
    \cdots \longrightarrow H^n(S.X_2) \longrightarrow H^n(S.X_3)
      \xrightarrow{\ \delta^n\ } H^{n+1}(S.X_1)
      \longrightarrow H^{n+1}(S.X_2) \longrightarrow \cdots .
  \]
  This is the complex-level homology long exact sequence and its connecting
  homomorphism. The derived-functor long exact sequence used in the sequel is
  constructed from this complex-level result via a horseshoe lift; see
  Lemma~\ref{mk-lem:acyclic_dimension_shift}.
\end{lemma}

\section{Right-acyclic objects}

\begin{definition}[Right-\(G\)-acyclic object]\label{mk-def:right_acyclic}
  \lean{CategoryTheory.Functor.IsRightAcyclic}
  \leanok
  \dcref{0157,custom}

  \uses{mk-lem:right_derived_injective_resolution}
  \textit{Source: Stacks Project, Derived Categories, Tag 0157
  (\texttt{definition-derived-functor}) and Tag 015C
  (\texttt{lemma-F-acyclic}).}
  Let \(G : \mathcal{A} \to \mathcal{B}\) be an additive functor between abelian
  categories. An object \(J\) of \(\mathcal{A}\) is \emph{right-\(G\)-acyclic} when
  the higher right-derived functors of \(G\) vanish at \(J\):
  \[
    (R^k G)(J) = 0 \qquad \text{for all } k \geq 1.
  \]
  Equivalently, writing the condition with the index shifted, \((R^{k+1} G)(J) = 0\)
  for all \(k \in \mathbb{N}\). The Stacks Project defines a right-\(G\)-acyclic
  object intrinsically (as one for which the one-term complex \(J[0]\) computes the
  derived functor \(RG\)); for an additive --- in our applications, left exact ---
  functor \(G\) that intrinsic condition is equivalent to the vanishing of all
  \((R^k G)(J)\) with \(k \geq 1\), which is the form we adopt here and which is
  directly checkable against the injective-resolution computation of
  Lemma~\ref{mk-lem:right_derived_injective_resolution}.

  \emph{Remark (alternative quantifier shapes).} Two equivalent formulations are
  available: ``\((R^k G)(J) = 0\) for all \(k \geq 1\)'' or the index-shifted
  ``\((R^{k+1} G)(J) = 0\) for all \(k \in \mathbb{N}\)''; the latter avoids the
  inequality side-condition and aligns with the vanishing statement of
  Lemma~\ref{mk-lem:right_derived_vanishes_injective}.
\end{definition}

\section{Lifting a short exact sequence to injective resolutions}

The passage from the complex-level homology long exact sequence
(Lemma~\ref{mk-lem:homology_long_exact_sequence}) to a long exact sequence of
right-derived functors requires one geometric input: a short exact sequence of
objects must be lifted to a short exact sequence of their injective resolutions
which is, moreover, \emph{degreewise split}. This is the dual Horseshoe Lemma. It is the single genuinely new homological-algebra construction the argument
requires; individual injective resolutions (provided by enough injectives) furnish
the building blocks, and the horseshoe construction assembles them into the required
lift.

\subsection{The horseshoe construction, step by step}

The construction proceeds by induction on degree; we isolate its independent moving
parts as the sub-lemmas below, from which the horseshoe lemma is assembled. Throughout,
\(I_A\) and \(I_C\) are fixed injective resolutions of \(A\) and \(C\) (they exist
because \(\mathcal{A}\) has enough injectives), and we write
\(d_A, d_C\) for their differentials and \(I_A^n, I_C^n\) for their degree-\(n\) terms.
The middle resolution is built degreewise as the biproduct
\(I_B^n := I_A^n \oplus I_C^n\). All seven statements are structural pieces of the dual
Horseshoe Lemma (dual of Weibel, \emph{An Introduction to Homological Algebra}, Lemma
2.2.8); no new source is introduced.

\begin{lemma}[Degree term: a finite biproduct of injectives is injective]
  \label{mk-lem:horseshoe_biprod_injective}
  \lean{CategoryTheory.Injective.instBiprod}
  \mathlibok
  \dcref{custom}

  In an abelian category, if \(P\) and \(Q\) are injective objects then their biproduct
  \(P \oplus Q\) is injective. Consequently each degree term
  \(I_B^n := I_A^n \oplus I_C^n\) of the middle complex is injective, being the
  biproduct of the injective terms \(I_A^n\) and \(I_C^n\) of the two chosen
  resolutions.
\end{lemma}

\begin{lemma}[Degree term: the biproduct short complex is split exact]
  \label{mk-lem:horseshoe_degree_split}
  \lean{CategoryTheory.ShortComplex.Splitting.ofHasBinaryBiproduct}
  \mathlibok
  \dcref{custom}

  For objects \(P, Q\) of an abelian category the short complex
  \[
    P \xrightarrow{\ \iota_P\ } P \oplus Q \xrightarrow{\ \pi_Q\ } Q
  \]
  formed from the biproduct coprojection \(\iota_P\) and projection \(\pi_Q\) carries a
  canonical \texttt{ShortComplex.Splitting}; in particular it is a split short exact
  sequence. Applied in each degree \(n\) with \(P = I_A^n\), \(Q = I_C^n\), this exhibits
  every degree
  \[
    0 \longrightarrow I_A^n \longrightarrow I_B^n \longrightarrow I_C^n \longrightarrow 0
  \]
  as split, furnishing the degreewise-splitting data required for the construction.
\end{lemma}

\begin{lemma}[Per-stage monomorphism into the biproduct]\label{mk-lem:horseshoe_stage_mono}
  \lean{CategoryTheory.mono_biprod_lift_factorThru_of_exact}
  \leanok
  \dcref{custom}

  \uses{mk-lem:horseshoe_biprod_injective}
  Let \(0 \to S.X_1 \xrightarrow{S.f} S.X_2 \xrightarrow{S.g} S.X_3 \to 0\) be an exact
  short complex with \(S.f\) a monomorphism, and let \(\alpha : S.X_1 \to P\),
  \(\gamma : S.X_3 \to Q\) be monomorphisms with \(P\) injective. Then the map
  \[
    \bigl\langle\, \mathrm{factorThru}(\alpha, S.f)\,,\; S.g \circ \gamma \,\bigr\rangle
      \;:\; S.X_2 \longrightarrow P \oplus Q,
  \]
  whose first component is the injective extension of \(\alpha\) along the mono \(S.f\)
  and whose second component is \(S.X_2 \twoheadrightarrow S.X_3 \hookrightarrow Q\), is a
  monomorphism. This is the kernel/cokernel step driving each stage of the recursion:
  applied to the \(n\)-th cosyzygy short exact sequence with \(P = I_A^{n+1}\),
  \(Q = I_C^{n+1}\) it produces the monomorphism of the next cosyzygy into the next
  biproduct term, whose cokernel feeds the following stage (and, at the base stage with
  the original sequence \(0 \to A \to B \to C \to 0\), the augmentation
  \(\beta : B \hookrightarrow I_B^0\)).
\end{lemma}

\begin{proof}
  \leanok
  \uses{mk-lem:horseshoe_biprod_injective}
  Write \(u := \langle \mathrm{factorThru}(\alpha, S.f), S.g \circ \gamma\rangle\). To see
  \(u\) is mono it suffices, in an abelian category, to show \(x \circ u = 0 \implies x = 0\)
  for any \(x\). Composing \(x \circ u\) with the two biproduct projections gives
  \(x \circ \mathrm{factorThru}(\alpha, S.f) = 0\) and \(x \circ S.g \circ \gamma = 0\);
  since \(\gamma\) is mono the latter yields \(x \circ S.g = 0\). Exactness of \(S\) at
  \(S.X_2\) (with \(S.f\) mono) lets us lift \(x\) through \(S.f\): there is \(y\) with
  \(y \circ S.f = x\). Then \(y \circ \alpha = y \circ S.f \circ \mathrm{factorThru}(\alpha, S.f)
  = x \circ \mathrm{factorThru}(\alpha, S.f) = 0\), using the defining property
  \(S.f \circ \mathrm{factorThru}(\alpha, S.f) = \alpha\); as \(\alpha\) is mono, \(y = 0\),
  hence \(x = y \circ S.f = 0\).
\end{proof}

\begin{lemma}[Off-diagonal twist and augmentation lifts]
  \label{mk-lem:horseshoe_twist}
  \lean{CategoryTheory.InjectiveResolution.horseshoeτ, CategoryTheory.InjectiveResolution.horseshoeτ_cocycle, CategoryTheory.InjectiveResolution.horseshoeβ, CategoryTheory.InjectiveResolution.twistPair,
    CategoryTheory.InjectiveResolution.horseshoeβ₁, CategoryTheory.InjectiveResolution.f_comp_horseshoeβ₁,
    CategoryTheory.InjectiveResolution.horseshoeH, CategoryTheory.InjectiveResolution.g_comp_horseshoeH,
    CategoryTheory.InjectiveResolution.horseshoeH_comp_d, CategoryTheory.InjectiveResolution.horseshoeτZero,
    CategoryTheory.InjectiveResolution.horseshoeτZero_hf, CategoryTheory.InjectiveResolution.horseshoeτ_zero,
    CategoryTheory.InjectiveResolution.horseshoeβ_comp_d, CategoryTheory.InjectiveResolution.horseshoeβ_fst,
    CategoryTheory.InjectiveResolution.horseshoeβ_snd, CategoryTheory.InjectiveResolution.ιC0,
    CategoryTheory.InjectiveResolution.ιC0_comp_d, CategoryTheory.InjectiveResolution.ιC0_comp_τZero,
    CategoryTheory.InjectiveResolution.ιC_comp_horseshoeτZero}
  \dcref{custom}

  \uses{mk-lem:horseshoe_stage_mono}
  There exist an augmentation \(\beta : B \to I_B^0\) and a family of morphisms
  \[
    \tau^n : I_C^n \longrightarrow I_A^{n+1} \qquad (n \in \mathbb{N})
  \]
  --- the off-diagonal components of the middle differential --- satisfying the cocycle
  identity
  \[
    \tau^{n+1} \circ d_C^n \;=\; -\, d_A^{n+1} \circ \tau^n \qquad (n \in \mathbb{N}),
  \]
  together with the compatibility of \(\beta\) with the augmentations \(A \to I_A^0\),
  \(C \to I_C^0\) and the maps of the original short exact sequence. Each \(\tau^n\) (and
  \(\beta\)) is produced by the universal lifting property of injectives: the target
  summand \(I_A^{n+1}\) is injective and the relevant source map is a monomorphism
  (Lemma~\ref{mk-lem:horseshoe_stage_mono}).
\end{lemma}

\begin{proof}
  \uses{mk-lem:horseshoe_stage_mono}
  The lifts are constructed by recursion on \(n\). For the augmentation, the first
  component of \(\beta\) extends the composite \(A \to I_A^0\) along the monomorphism
  \(A \hookrightarrow B\) using injectivity of \(I_A^0\); the second component is
  \(B \twoheadrightarrow C \to I_C^0\).
  At the inductive step we have \(\tau^{n}\) (with \(\tau^{n} : I_C^{n} \to I_A^{n+1}\))
  and seek \(\tau^{n+1} : I_C^{n+1} \to I_A^{n+2}\) with
  \(\tau^{n+1} \circ d_C^n = -\,d_A^{n+1} \circ \tau^n\). The right-hand side
  \(-\,d_A^{n+1} \circ \tau^n : I_C^n \to I_A^{n+2}\) annihilates
  \(\ker d_C^n = \operatorname{im} d_C^{n-1}\) --- the equality \(\ker d_C^n =
  \operatorname{im} d_C^{n-1}\) being the exactness of the injective resolution \(I_C\) in
  the positive degree \(n\), which is part of the defining data of an injective resolution.
  Indeed, precomposing with \(d_C^{n-1}\) and using the
  previous-stage identity \(\tau^{n} \circ d_C^{n-1} = -\,d_A^{n} \circ \tau^{n-1}\) gives
  \(-\,d_A^{n+1} \circ \tau^n \circ d_C^{n-1} = d_A^{n+1} \circ d_A^{n} \circ \tau^{n-1} = 0\),
  where the final equality is the complex identity \(d_A^{n+1} \circ d_A^{n} = 0\) for the
  cochain complex \(I_A\). Hence
  \(-\,d_A^{n+1} \circ \tau^n\) factors through the image of \(d_C^n\), i.e.\ through the
  monomorphism \(\operatorname{im} d_C^n \hookrightarrow I_C^{n+1}\); injectivity of the
  target \(I_A^{n+2}\) extends that factorization along this monomorphism to a map
  \(\tau^{n+1} : I_C^{n+1} \to I_A^{n+2}\) defined on all of \(I_C^{n+1}\) and satisfying
  \(\tau^{n+1} \circ d_C^n = -\,d_A^{n+1} \circ \tau^n\) by construction. This is exactly
  the dual of the inductive lift in the Horseshoe Lemma.
\end{proof}

\begin{lemma}[The assembled middle differential squares to zero]\label{mk-lem:horseshoe_dComp}
  \lean{CategoryTheory.twistedBiprodD_comp}
  \leanok
  \dcref{custom}

  \uses{mk-lem:horseshoe_twist}
  With the twist family \(\tau\) of Lemma~\ref{mk-lem:horseshoe_twist}, define the middle
  differential degreewise by the matrix
  \[
    d_B^n \;=\;
    \begin{pmatrix} d_A^n & \tau^n \\ 0 & d_C^n \end{pmatrix}
    \;:\; I_A^n \oplus I_C^n \longrightarrow I_A^{n+1} \oplus I_C^{n+1}.
  \]
  Then \(d_B^{n+1} \circ d_B^n = 0\) for every \(n\), so \((I_B^\bullet, d_B)\) is a
  cochain complex.
\end{lemma}

\begin{proof}
  \leanok
  \uses{mk-lem:horseshoe_twist}
  Multiplying the two matrices, the composite \(d_B^{n+1} \circ d_B^n\) has diagonal
  entries \(d_A^{n+1} \circ d_A^n = 0\) and \(d_C^{n+1} \circ d_C^n = 0\) (the complexes
  \(I_A\), \(I_C\) being complexes), lower-left entry \(0\), and upper-right entry
  \[
    d_A^{n+1} \circ \tau^n + \tau^{n+1} \circ d_C^n,
  \]
  which vanishes precisely by the cocycle identity
  \(\tau^{n+1} \circ d_C^n = -\,d_A^{n+1} \circ \tau^n\) of
  Lemma~\ref{mk-lem:horseshoe_twist}. Hence \(d_B^{n+1} \circ d_B^n = 0\).
\end{proof}

\begin{lemma}[The inclusion and projection are chain maps]
  \label{mk-lem:horseshoe_chainMap}
  \lean{CategoryTheory.twistedBiprodInl, CategoryTheory.twistedBiprodSnd, CategoryTheory.twistedBiprodSplitting, CategoryTheory.InjectiveResolution.horseshoeSES_shortExact,
    CategoryTheory.InjectiveResolution.horseshoeMid, CategoryTheory.InjectiveResolution.horseshoeSES,
    CategoryTheory.InjectiveResolution.horseshoeSES_splitting, CategoryTheory.InjectiveResolution.horseshoeι,
    CategoryTheory.InjectiveResolution.horseshoeι_f_zero, CategoryTheory.InjectiveResolution.mono_horseshoeβ,
    CategoryTheory.twistedBiprod, CategoryTheory.twistedBiprod_X, CategoryTheory.twistedBiprod_d,
    CategoryTheory.twistedBiprodD, CategoryTheory.twistedBiprodD_fst, CategoryTheory.twistedBiprodD_snd,
    CategoryTheory.twistedBiprodInl_f, CategoryTheory.twistedBiprodSnd_f, CategoryTheory.twistedBiprodInl_comp_Snd}
  \dcref{custom}

  \uses{mk-lem:horseshoe_dComp}
  With \(d_B\) as in Lemma~\ref{mk-lem:horseshoe_dComp}, the degreewise biproduct
  coprojection \(\iota^n = \iota_{I_A^n} : I_A^n \to I_B^n\) and projection
  \(\pi^n = \pi_{I_C^n} : I_B^n \to I_C^n\) assemble into chain maps
  \(\iota : I_A^\bullet \to I_B^\bullet\) and \(\pi : I_B^\bullet \to I_C^\bullet\); that
  is, \(d_B^n \circ \iota^n = \iota^{n+1} \circ d_A^n\) and
  \(\pi^{n+1} \circ d_B^n = d_C^n \circ \pi^n\). Moreover \(\iota, \pi\) are compatible
  with the augmentations and form, in each degree, the split short exact sequence
  \(0 \to I_A^n \to I_B^n \to I_C^n \to 0\) of
  Lemma~\ref{mk-lem:horseshoe_degree_split}.
\end{lemma}

\begin{proof}
  \uses{mk-lem:horseshoe_dComp, mk-lem:horseshoe_degree_split}
  These are direct biproduct computations from the matrix form of \(d_B\). Pre-composing
  \(d_B^n\) with the coprojection \(\iota^n\) of the first summand selects the first
  column \(\begin{pmatrix} d_A^n \\ 0\end{pmatrix}\), which equals
  \(\iota^{n+1} \circ d_A^n\); so \(\iota\) is a chain map. Post-composing \(d_B^n\) with
  the projection \(\pi^{n+1}\) onto the second summand selects the second row
  \(\begin{pmatrix} 0 & d_C^n\end{pmatrix}\), which equals \(d_C^n \circ \pi^n\); so
  \(\pi\) is a chain map. Compatibility with the augmentations is the corresponding
  statement in degree \(-1 \to 0\), immediate from the construction of \(\beta\) in
  Lemma~\ref{mk-lem:horseshoe_twist}. In each degree \(\iota^n = \iota_{I_A^n}\) and
  \(\pi^n = \pi_{I_C^n}\) are the structure maps of the biproduct, so the degree is split
  exact by Lemma~\ref{mk-lem:horseshoe_degree_split}.
\end{proof}

\begin{lemma}[Middle-term quasi-isomorphism transfer]\label{mk-lem:quasiIso_tau2}
  \lean{HomologicalComplex.HomologySequence.quasiIso_τ₂}
  \leanok
  \dcref{custom}

  \uses{mk-lem:homology_long_exact_sequence}
  Project-local supplement (the \(\tau_2\) companion of Mathlib's \(\tau_3\) transfer
  \texttt{quasiIso\_τ₃}). Let \(\varphi : S_1 \to S_2\) be a morphism of short exact
  sequences of cochain complexes, with components \(\varphi.\tau_1, \varphi.\tau_2,
  \varphi.\tau_3\) on the three columns. Suppose \(\varphi.\tau_1\) and \(\varphi.\tau_3\)
  are quasi-isomorphisms, and that the middle component \(\varphi.\tau_2\) is a
  monomorphism in every degree \(i\) that has no incoming differential, and an
  epimorphism in every degree \(i\) that has no outgoing differential. Then
  \(\varphi.\tau_2\) is a quasi-isomorphism.
\end{lemma}

\begin{proof}
  \leanok
  \uses{mk-lem:homology_long_exact_sequence}
  This is the homology four-lemma. Fix a degree \(i\); we show
  \(H^i(\varphi.\tau_2)\) is an isomorphism by proving it is both epimorphic and
  monomorphic and concluding that a map which is both is an isomorphism. \emph{Epi:}
  when \(i\) has an outgoing differential to some \(j\), the two short exact sequences
  yield a commuting ladder whose rows are the five-term exact windows of the homology
  long exact sequence (Lemma~\ref{mk-lem:homology_long_exact_sequence}) about degree
  \(i\); the four-lemma (epi form), using that \(H^i(\varphi.\tau_1)\),
  \(H^i(\varphi.\tau_3)\) are epi and \(H^j(\varphi.\tau_1)\) is mono, forces
  \(H^i(\varphi.\tau_2)\) epi. When \(i\) has no outgoing differential,
  \(\varphi.\tau_2\) is epi in degree \(i\) by hypothesis, and an epimorphism of
  complexes that is epi in a degree with no outgoing differential is epi on homology
  there. \emph{Mono:} dually, when \(i\) has an incoming differential from some \(k\),
  the four-lemma (mono form), using that \(H^k(\varphi.\tau_3)\) is epi and
  \(H^i(\varphi.\tau_1), H^i(\varphi.\tau_3)\) are mono, forces
  \(H^i(\varphi.\tau_2)\) mono; when \(i\) has no incoming differential it is mono in
  that degree by hypothesis, hence mono on homology. Being epi and mono in every
  degree, \(H^i(\varphi.\tau_2)\) is an isomorphism for all \(i\), so
  \(\varphi.\tau_2\) is a quasi-isomorphism.
\end{proof}

\begin{lemma}[The middle complex resolves \(B\)]
  \label{mk-lem:horseshoe_resolvesMiddle}
  \lean{CategoryTheory.InjectiveResolution.ofShortExact_resolvesMiddle,
    CategoryTheory.InjectiveResolution.horseshoeφ, CategoryTheory.InjectiveResolution.horseshoeφ_comm₁₂,
    CategoryTheory.InjectiveResolution.horseshoeφ_comm₂₃, CategoryTheory.InjectiveResolution.horseshoeφ_τ₁,
    CategoryTheory.InjectiveResolution.horseshoeφ_τ₂, CategoryTheory.InjectiveResolution.horseshoeφ_τ₃,
    CategoryTheory.InjectiveResolution.quasiIso_horseshoeι, CategoryTheory.InjectiveResolution.single₀_hom_ext}
  \dcref{custom}

  \uses{mk-lem:horseshoe_chainMap, mk-lem:homology_long_exact_sequence, mk-lem:quasiIso_tau2}
  The cochain complex \((I_B^\bullet, d_B)\) of Lemma~\ref{mk-lem:horseshoe_dComp}, with the
  augmentation \(\beta : B \to I_B^0\) of Lemma~\ref{mk-lem:horseshoe_twist}, is an injective
  resolution of \(B\): its terms are injective
  (Lemma~\ref{mk-lem:horseshoe_biprod_injective}), it is exact in every positive degree, and
  \(H^0 \cong B\) via \(\beta\).
\end{lemma}

\begin{proof}
  \uses{mk-lem:horseshoe_chainMap, mk-lem:homology_long_exact_sequence,
    mk-lem:horseshoe_biprod_injective, mk-lem:quasiIso_tau2}
  By Lemma~\ref{mk-lem:horseshoe_chainMap} the maps \(\iota, \pi\) form, degreewise, the
  split short exact sequence \(0 \to I_A^n \to I_B^n \to I_C^n \to 0\); assembling these
  splittings (Lemma~\ref{mk-lem:horseshoe_degree_split}) gives a short exact sequence of
  cochain complexes
  \[
    0 \longrightarrow I_A^\bullet \longrightarrow I_B^\bullet
      \longrightarrow I_C^\bullet \longrightarrow 0 .
  \]
  Apply the complex-level homology long exact sequence
  (Lemma~\ref{mk-lem:homology_long_exact_sequence}). The outer complexes \(I_A^\bullet\),
  \(I_C^\bullet\) are injective resolutions, hence exact in every positive degree
  (\(H^{k}(I_A) = H^{k}(I_C) = 0\) for \(k \geq 1\)) with \(H^0(I_A) \cong A\),
  \(H^0(I_C) \cong C\). In the long exact sequence the terms flanking \(H^{k}(I_B)\) for
  \(k \geq 1\) vanish, forcing \(H^{k}(I_B) = 0\); in degree \(0\) the initial segment
  \(0 \to H^0(I_A) \to H^0(I_B) \to H^0(I_C) \to H^1(I_A) = 0\) reads
  \(0 \to A \to H^0(I_B) \to C \to 0\), which together with the original sequence
  \(0 \to A \to B \to C \to 0\) and the augmentation \(\beta\) identifies
  \(H^0(I_B) \cong B\). Hence \(I_B^\bullet\) is exact in positive degrees and resolves
  \(B\); with injective terms it is an injective resolution of \(B\).
\end{proof}

\begin{lemma}[Horseshoe lift of a short exact sequence to injective resolutions]\label{mk-lem:injective_resolution_of_ses}
  \lean{CategoryTheory.InjectiveResolution.ofShortExact}
  \leanok
  \dcref{custom}

  \uses{mk-lem:horseshoe_biprod_injective, mk-lem:horseshoe_degree_split,
    mk-lem:horseshoe_stage_mono, mk-lem:horseshoe_twist, mk-lem:horseshoe_dComp,
    mk-lem:horseshoe_chainMap, mk-lem:horseshoe_resolvesMiddle,
    mk-lem:right_derived_injective_resolution}
  Let \(\mathcal{A}\) be an abelian category with enough injectives and let
  \[
    0 \longrightarrow A \longrightarrow B \longrightarrow C \longrightarrow 0
  \]
  be a short exact sequence in \(\mathcal{A}\). Then there exist injective
  resolutions \(I_A\) of \(A\), \(I_B\) of \(B\), \(I_C\) of \(C\), together with a
  short exact sequence of cochain complexes
  \[
    0 \longrightarrow I_A \longrightarrow I_B \longrightarrow I_C \longrightarrow 0
  \]
  which is compatible with the augmentations \(A \to I_A^0\), \(B \to I_B^0\),
  \(C \to I_C^0\) and which is \emph{degreewise split}: for every \(n\) the sequence
  \[
    0 \longrightarrow I_A^n \longrightarrow I_B^n \longrightarrow I_C^n
      \longrightarrow 0
  \]
  is a split short exact sequence of \(\mathcal{A}\), so that
  \(I_B^n \cong I_A^n \oplus I_C^n\).
\end{lemma}

\begin{proof}
  \leanok
  \uses{mk-lem:horseshoe_biprod_injective, mk-lem:horseshoe_degree_split,
    mk-lem:horseshoe_stage_mono, mk-lem:horseshoe_twist, mk-lem:horseshoe_dComp,
    mk-lem:horseshoe_chainMap, mk-lem:horseshoe_resolvesMiddle,
    mk-lem:right_derived_injective_resolution, mk-lem:homology_long_exact_sequence}
  This is the (dual) Horseshoe Lemma (dual of Weibel, \emph{An Introduction to
  Homological Algebra}, Lemma 2.2.8), assembled from the sub-lemmas of \S``The horseshoe
  construction, step by step''. Take \(I_A, I_C\) and set \(I_B^n := I_A^n \oplus I_C^n\):
  injective by Lemma~\ref{mk-lem:horseshoe_biprod_injective}, degreewise split by
  Lemma~\ref{mk-lem:horseshoe_degree_split}. The twist \(\tau\) and augmentation exist with
  the cocycle identity (Lemma~\ref{mk-lem:horseshoe_twist}, via
  Lemma~\ref{mk-lem:horseshoe_stage_mono}); the matrix differential then squares to zero
  (Lemma~\ref{mk-lem:horseshoe_dComp}) with coprojection and projection chain maps
  (Lemma~\ref{mk-lem:horseshoe_chainMap}); and \(I_B^\bullet\) resolves \(B\)
  (Lemma~\ref{mk-lem:horseshoe_resolvesMiddle}). These data are exactly the asserted
  degreewise-split short exact sequence \(0 \to I_A \to I_B \to I_C \to 0\) of cochain
  complexes compatible with the augmentations.
\end{proof}

\section{The dimension-shift step}

The induction driving the comparison theorem rests on a single short-exact-sequence
calculation: cutting a resolution into short exact sequences and reading off the
long exact sequence of right-derived functors. The latter is not available off the
shelf at the level of the derived-functor objects \((R^k G)\); we build it from the
complex-level homology long exact sequence
(Lemma~\ref{mk-lem:homology_long_exact_sequence}) via the horseshoe lift
(Lemma~\ref{mk-lem:injective_resolution_of_ses}). We first record the shift at the
\emph{resolution} level — where the horseshoe lift directly supplies the
degreewise-split short exact sequence of injective resolutions — and then specialize
to objects.

\begin{lemma}[Resolution-level dimension shift]
  \label{mk-lem:right_derived_shift_split_resolution}
  \lean{CategoryTheory.Functor.rightDerivedShiftIsoOfSplitResolutionSES, CategoryTheory.shortExact_of_degreewise_splitting, CategoryTheory.shortExact_map_mapHomologicalComplex_of_degreewise_splitting, CategoryTheory.Functor.isZero_homology_mapHomologicalComplex_of_isRightAcyclic}
  \dcref{custom}

  \uses{mk-def:right_acyclic, mk-lem:right_derived_injective_resolution,
    mk-lem:homology_long_exact_sequence}
  Let \(G : \mathcal{A} \to \mathcal{B}\) be additive and let
  \(0 \to A \to J \to Z \to 0\) be a short exact sequence with \(J\) right-\(G\)-acyclic,
  lifted to a \emph{degreewise-split} short exact sequence of injective resolutions
  \(0 \to I_A \to I_J \to I_Z \to 0\) (chain maps \(\varphi, \psi\) with
  \(\psi \circ \varphi = 0\) and a splitting of \(0 \to I_A^n \to I_J^n \to I_Z^n \to 0\)
  in each degree \(n\)). Then for every \(k \geq 0\) the connecting map of the homology
  long exact sequence of \(G(I_\bullet)\) is an isomorphism
  \[
    (R^{k+1} G)(Z) \;\cong\; (R^{k+2} G)(A).
  \]
\end{lemma}

\begin{proof}
  \uses{mk-def:right_acyclic, mk-lem:right_derived_injective_resolution,
    mk-lem:homology_long_exact_sequence}
  Applying \(G\) degreewise to the degreewise-split sequence preserves each split
  degree, yielding a short exact sequence of cochain complexes
  \(0 \to G(I_A) \to G(I_J) \to G(I_Z) \to 0\) in \(\mathcal{B}\). Its homology long
  exact sequence (Lemma~\ref{mk-lem:homology_long_exact_sequence}) has connecting maps
  \(\delta^m : H^m(G(I_Z)) \to H^{m+1}(G(I_A))\); transport these across the
  injective-resolution computation \(H^m(G(I_-)) \cong (R^m G)(-)\)
  (Lemma~\ref{mk-lem:right_derived_injective_resolution}). For \(m = k+1 \geq 1\) the two
  homology objects flanking \(\delta^m\) in the long exact sequence both sit on the
  middle complex \(G(I_J)\), namely \(H^{k+1}(G(I_J))\) and \(H^{k+2}(G(I_J))\); both
  vanish because \(J\) is right-\(G\)-acyclic (Definition~\ref{mk-def:right_acyclic}).
  With both neighbours zero, \(\delta^{k+1}\) is an isomorphism, giving
  \((R^{k+1} G)(Z) \cong (R^{k+2} G)(A)\).
\end{proof}

\begin{lemma}[Dimension shift across an acyclic short exact sequence]\label{mk-lem:acyclic_dimension_shift}
  \lean{CategoryTheory.Functor.rightDerivedShiftIsoOfAcyclic}
  \leanok
  \dcref{015D,custom}

  \uses{mk-def:right_acyclic, mk-lem:right_derived_zero_iso_self,
    mk-lem:injective_resolution_of_ses, mk-lem:homology_long_exact_sequence,
    mk-lem:right_derived_shift_split_resolution}
  \textit{Source: Stacks Project, Derived Categories, Tag 015D
  (\texttt{lemma-F-acyclic-ses}).}
  Let \(G : \mathcal{A} \to \mathcal{B}\) be an additive left-exact functor, and let
  \[
    0 \longrightarrow A \longrightarrow J \longrightarrow Z \longrightarrow 0
  \]
  be a short exact sequence in \(\mathcal{A}\) with the middle term \(J\)
  right-\(G\)-acyclic. Then the connecting maps of the long exact sequence of
  right-derived functors furnish isomorphisms
  \[
    (R^{k} G)(Z) \;\xrightarrow{\ \cong\ }\; (R^{k+1} G)(A)
    \qquad \text{for all } k \geq 1.
  \]
  (The companion lowest-degree statement
  \((R^{1} G)(A) \cong \operatorname{coker}(G(J) \to G(Z))\) is isolated as the
  sibling Lemma~\ref{mk-lem:acyclic_one_iso_coker}.)
\end{lemma}
\begin{proof}
  \leanok
  \uses{mk-def:right_acyclic, mk-lem:right_derived_zero_iso_self,
    mk-lem:right_derived_injective_resolution, mk-lem:homology_long_exact_sequence,
    mk-lem:injective_resolution_of_ses, mk-lem:right_derived_shift_split_resolution}
    We first construct the long exact sequence of right-derived functors associated
  to \(0 \to A \to J \to Z \to 0\). In the classical treatment this sequence is
  supplied by the \(\delta\)-functor structure of \(\{R^k G\}\) on the derived
  category; here it is built directly from the complex-level homology long exact
  sequence via a horseshoe lift.

  \emph{Building the long exact sequence.} By the horseshoe lift
  (Lemma~\ref{mk-lem:injective_resolution_of_ses}) choose injective resolutions
  \(I_A, I_J, I_Z\) of \(A, J, Z\) fitting into a short exact sequence of cochain
  complexes
  \[
    0 \longrightarrow I_A \longrightarrow I_J \longrightarrow I_Z \longrightarrow 0
  \]
  that is degreewise split. Apply \(G\) degreewise. Because each degree
  \(0 \to I_A^n \to I_J^n \to I_Z^n \to 0\) is a \emph{split} short exact sequence
  (so \(I_J^n \cong I_A^n \oplus I_Z^n\) compatibly with the maps), the additive
  functor \(G\) carries it to a split short exact sequence
  \(0 \to G(I_A^n) \to G(I_J^n) \to G(I_Z^n) \to 0\) --- this is where degreewise
  splitness is essential, since \(G\) is not assumed exact and would not in general
  preserve a non-split short exact sequence. Hence
  \[
    0 \longrightarrow G(I_A) \longrightarrow G(I_J) \longrightarrow G(I_Z)
      \longrightarrow 0
  \]
  is a short exact sequence of cochain complexes in \(\mathcal{B}\). The
  complex-level homology long exact sequence
  (Lemma~\ref{mk-lem:homology_long_exact_sequence}) applied to it reads
  \[
    \cdots \to H^k\bigl(G(I_J)\bigr) \to H^k\bigl(G(I_Z)\bigr)
      \xrightarrow{\ \delta^k\ } H^{k+1}\bigl(G(I_A)\bigr)
      \to H^{k+1}\bigl(G(I_J)\bigr) \to \cdots .
  \]
  By Lemma~\ref{mk-lem:right_derived_injective_resolution} the homology of each applied
  resolution computes the right-derived functor:
  \(H^k(G(I_A)) \cong (R^k G)(A)\), \(H^k(G(I_J)) \cong (R^k G)(J)\) and
  \(H^k(G(I_Z)) \cong (R^k G)(Z)\), naturally in the maps of the sequence. Transporting
  the homology long exact sequence across these isomorphisms gives the long exact
  sequence of right-derived functors
  \[
    \cdots \to (R^k G)(J) \to (R^k G)(Z)
      \xrightarrow{\ \delta^k\ } (R^{k+1} G)(A) \to (R^{k+1} G)(J) \to \cdots,
  \]
  and, using \(R^0 G \cong G\) (Lemma~\ref{mk-lem:right_derived_zero_iso_self}) at the
  bottom, its low-degree initial segment
  \[
    0 \to G(A) \to G(J) \to G(Z)
      \xrightarrow{\ \delta^0\ } (R^1 G)(A) \to (R^1 G)(J) \to \cdots .
  \]

  \emph{Killing the acyclic terms.} Because \(J\) is right-\(G\)-acyclic
  (Definition~\ref{mk-def:right_acyclic}), every term \((R^k G)(J)\) with \(k \geq 1\)
  vanishes. For \(k \geq 1\) the segment around \((R^k G)(J)\),
  \[
    \underbrace{(R^{k} G)(J)}_{=0} \to (R^{k} G)(Z)
      \xrightarrow{\ \delta^k\ } (R^{k+1} G)(A) \to \underbrace{(R^{k+1} G)(J)}_{=0},
  \]
  has both outer terms zero, so the connecting map
  \(\delta^k : (R^{k} G)(Z) \to (R^{k+1} G)(A)\) is an isomorphism, which is the
  asserted dimension shift. (The lowest-degree segment of the same long exact
  sequence yields the cokernel description of \((R^1 G)(A)\); that reading is
  recorded separately in Lemma~\ref{mk-lem:acyclic_one_iso_coker}.)
\end{proof}

\begin{lemma}[Lowest-degree cokernel description across an acyclic short exact sequence]\label{mk-lem:acyclic_one_iso_coker}
  \lean{CategoryTheory.Functor.rightDerivedOneIsoCokerOfAcyclic}
  \leanok
  \dcref{015D,custom}

  \uses{mk-def:right_acyclic, mk-lem:right_derived_zero_iso_self,
    mk-lem:right_derived_injective_resolution, mk-lem:homology_long_exact_sequence,
    mk-lem:injective_resolution_of_ses}
  \textit{Source: Stacks Project, Derived Categories, Tag 015D
  (\texttt{lemma-F-acyclic-ses}).}
  Let \(G : \mathcal{A} \to \mathcal{B}\) be an additive left-exact functor, and let
  \[
    0 \longrightarrow A \longrightarrow J \longrightarrow Z \longrightarrow 0
  \]
  be a short exact sequence in \(\mathcal{A}\) with the middle term \(J\)
  right-\(G\)-acyclic. Then there is an isomorphism in \(\mathcal{B}\)
  \[
    (R^{1} G)(A) \;\cong\; \operatorname{coker}\bigl(G(J) \to G(Z)\bigr),
  \]
  i.e.\ \((R^1 G)(A)\) is the cokernel of the map \(G(J) \to G(Z)\) induced by the
  surjection \(J \twoheadrightarrow Z\).
\end{lemma}

\begin{proof}
  \leanok
  \uses{mk-def:right_acyclic, mk-lem:right_derived_zero_iso_self,
    mk-lem:right_derived_injective_resolution, mk-lem:homology_long_exact_sequence,
    mk-lem:injective_resolution_of_ses}
    This reads off the lowest-degree segment of the very long exact sequence of
  right-derived functors built, from the same short exact sequence
  \(0 \to A \to J \to Z \to 0\), in the proof of
  Lemma~\ref{mk-lem:acyclic_dimension_shift}; here we use only its initial terms, so we
  recall the construction briefly and then specialize. By the horseshoe lift
  (Lemma~\ref{mk-lem:injective_resolution_of_ses}) choose injective resolutions
  \(I_A, I_J, I_Z\) of \(A, J, Z\) fitting into a degreewise-split short exact
  sequence of cochain complexes \(0 \to I_A \to I_J \to I_Z \to 0\). Applying the
  additive functor \(G\) degreewise preserves each split degree, giving a short exact
  sequence of complexes \(0 \to G(I_A) \to G(I_J) \to G(I_Z) \to 0\) in
  \(\mathcal{B}\); its complex-level homology long exact sequence
  (Lemma~\ref{mk-lem:homology_long_exact_sequence}), transported across the
  injective-resolution computation \(H^k(G(I_\bullet)) \cong (R^k G)(-)\)
  (Lemma~\ref{mk-lem:right_derived_injective_resolution}) and using
  \(R^0 G \cong G\) (Lemma~\ref{mk-lem:right_derived_zero_iso_self}), has low-degree
  initial segment
  \[
    G(J) \to G(Z) \xrightarrow{\ \delta^0\ } (R^1 G)(A) \to (R^1 G)(J).
  \]
  Because \(J\) is right-\(G\)-acyclic (Definition~\ref{mk-def:right_acyclic}), the right
  end \((R^1 G)(J)\) vanishes. Hence the segment
  \(G(J) \to G(Z) \xrightarrow{\ \delta^0\ } (R^1 G)(A) \to 0\) is exact with vanishing
  right end, so \(\delta^0\) is an epimorphism whose kernel is the image of
  \(G(J) \to G(Z)\). Therefore \(\delta^0\) factors through an isomorphism
  \(\operatorname{coker}\bigl(G(J) \to G(Z)\bigr) \xrightarrow{\ \cong\ } (R^1 G)(A)\),
  which is the assertion.
\end{proof}

\section{Cosyzygies of an acyclic resolution}

The comparison theorem decomposes an acyclic resolution into the short exact
sequences cut out by its cosyzygies, and reads the cohomology of the applied complex
off those sequences using left-exactness of \(G\). We isolate the two ingredients
here.

\begin{lemma}[Cosyzygy short exact sequences of an acyclic resolution]
  \label{mk-lem:cosyzygy_ses}
  \lean{CategoryTheory.Functor.cosyzygyShortExact, CategoryTheory.Functor.cosyzygyShortComplex, CategoryTheory.cosyzygyKernelFork, CategoryTheory.cosyzygy_iCycles_comp_toCycles, CategoryTheory.epi_toCycles_of_exactAt}
  \dcref{05TA,custom}

  \uses{mk-def:right_acyclic}
  \textit{Source: Stacks Project, Derived Categories, Tag 05TA
  (\texttt{proposition-enough-acyclics}).}
  Let \(G : \mathcal{A} \to \mathcal{B}\) be additive and let
  \[
    0 \longrightarrow A
      \longrightarrow J^0 \longrightarrow J^1 \longrightarrow J^2
      \longrightarrow \cdots
  \]
  be an exact resolution of \(A\) with every \(J^n\) right-\(G\)-acyclic. Define the
  cosyzygies \(Z^0 := A\) and, for \(m \geq 1\),
  \[
    Z^m \;:=\; \ker\bigl(J^m \to J^{m+1}\bigr)
      \;=\; \operatorname{im}\bigl(J^{m-1} \to J^m\bigr).
  \]
  Then for every \(m \geq 0\) the sequence
  \[
    0 \longrightarrow Z^m \longrightarrow J^m \longrightarrow Z^{m+1}
      \longrightarrow 0 \tag{$S_m$}
  \]
  is short exact, with the inclusion \(Z^m \hookrightarrow J^m\) the canonical one,
  the surjection \(J^m \twoheadrightarrow Z^{m+1}\) the corestriction of
  \(J^m \to J^{m+1}\) onto its image, and the middle term \(J^m\) right-\(G\)-acyclic.
  In particular \(Z^1 = \operatorname{im}(J^0 \to J^1) = \operatorname{coker}(A \to J^0)\).
\end{lemma}

\begin{proof}
  \uses{mk-def:right_acyclic}
  Fix \(m \geq 0\). Mono of \(Z^m \hookrightarrow J^m\) is immediate: \(Z^m\) is by
  definition a subobject of \(J^m\) (for \(m = 0\) via the injection \(A \hookrightarrow
  J^0\) of the resolution, for \(m \geq 1\) as the kernel \(\ker(J^m \to J^{m+1})\)).
  The map \(J^m \to Z^{m+1}\) is the corestriction of the differential
  \(d^m : J^m \to J^{m+1}\) onto its image \(Z^{m+1} = \operatorname{im}(d^m)\), hence
  an epimorphism. Exactness in the middle is exactness of the resolution at \(J^m\):
  the image of \(Z^m \hookrightarrow J^m\) equals \(\ker(d^m)\) --- for \(m \geq 1\)
  by the definition \(Z^m = \ker(J^m \to J^{m+1})\), and for \(m = 0\) because
  exactness of \(0 \to A \to J^0 \to J^1\) makes \(A \hookrightarrow J^0\) identify
  \(A\) with \(\ker(d^0)\) --- while the kernel of the epimorphism
  \(J^m \twoheadrightarrow Z^{m+1}\) is exactly \(\ker(d^m)\). Thus
  \(\operatorname{im}(Z^m \to J^m) = \ker(J^m \to Z^{m+1})\), so \((S_m)\) is short
  exact. The middle term \(J^m\) is right-\(G\)-acyclic by hypothesis. The final
  identity \(Z^1 = \operatorname{im}(J^0 \to J^1) = \operatorname{coker}(A \to J^0)\)
  is the case \(m = 0\): \(\operatorname{coker}(A \to J^0) = \operatorname{coker}(Z^0 \to J^0)
  \cong Z^1\) by \((S_0)\).
\end{proof}

\begin{lemma}[Applied cosyzygy is the cocycle object of the applied complex]
  \label{mk-lem:applied_cosyzygy_cycles}
  \lean{CategoryTheory.Functor.gCosyzygyIsoCocycles, CategoryTheory.Functor.gCosyzygyIsoCocycles_hom_iCycles, CategoryTheory.Functor.gCosyzygyIsoCocycles_toCycles}
  \dcref{015D,custom}

  \uses{mk-lem:cosyzygy_ses, mk-lem:right_derived_zero_iso_self}
  \textit{Source: Stacks Project, Derived Categories, Tag 015D
  (\texttt{lemma-F-acyclic-ses}), left-exact initial segment.}
  With \(G\) additive left exact and the acyclic resolution and cosyzygies of
  Lemma~\ref{mk-lem:cosyzygy_ses}, for every \(n \geq 0\) the morphism \(G\) applies to
  the inclusion \(Z^n \hookrightarrow J^n\) to give an isomorphism
  \[
    G(Z^n) \;\xrightarrow{\ \cong\ }\; \ker\bigl(G(J^n) \to G(J^{n+1})\bigr),
  \]
  identifying \(G(Z^n)\), as a subobject of \(G(J^n)\), with the \(n\)-th cocycle
  object of the applied complex \(G(J^\bullet)\). For \(n = 0\) this reads
  \(G(A) \cong \ker(G(J^0) \to G(J^1)) = H^0\bigl(G(J^\bullet)\bigr)\).
\end{lemma}

\begin{proof}
  \uses{mk-lem:cosyzygy_ses, mk-lem:right_derived_zero_iso_self}
  Fix \(n \geq 0\). Apply the left-exact functor \(G\) to the cosyzygy short exact
  sequence \((S_n)\), \(0 \to Z^n \to J^n \to Z^{n+1} \to 0\) of
  Lemma~\ref{mk-lem:cosyzygy_ses}; left-exactness yields an exact sequence
  \[
    0 \longrightarrow G(Z^n) \longrightarrow G(J^n) \longrightarrow G(Z^{n+1}),
  \]
  so \(G(Z^n) \hookrightarrow G(J^n)\) is the kernel of \(G(J^n) \to G(Z^{n+1})\).
  Next, the differential \(G(J^n) \to G(J^{n+1})\) of the applied complex factors as
  \(G(J^n) \xrightarrow{G(J^n \twoheadrightarrow Z^{n+1})} G(Z^{n+1})
  \xrightarrow{G(Z^{n+1} \hookrightarrow J^{n+1})} G(J^{n+1})\), because the original
  differential \(J^n \to J^{n+1}\) factors through its image \(Z^{n+1}\) as
  \(J^n \twoheadrightarrow Z^{n+1} \hookrightarrow J^{n+1}\) and \(G\) preserves
  composition. The second factor \(G(Z^{n+1}) \to G(J^{n+1})\) is a monomorphism (the
  image of the left-exact application of \(G\) to \((S_{n+1})\), whose first map
  \(G(Z^{n+1}) \hookrightarrow G(J^{n+1})\) is mono). A monomorphism does not change
  kernels under precomposition, so
  \(\ker(G(J^n) \to G(J^{n+1})) = \ker(G(J^n) \to G(Z^{n+1})) = G(Z^n)\). Hence
  \(G(Z^n) \xrightarrow{\cong} \ker(G(J^n) \to G(J^{n+1}))\) as subobjects of
  \(G(J^n)\). For \(n = 0\), \(Z^0 = A\) gives
  \(G(A) \cong \ker(G(J^0) \to G(J^1)) = H^0(G(J^\bullet))\), the last equality being
  the definition of the zeroth cohomology of a cochain complex starting in degree
  \(0\) (no incoming differential). The use of \(R^0 G \cong G\)
  (Lemma~\ref{mk-lem:right_derived_zero_iso_self}) is what licenses reading \(G(A)\) as
  \((R^0 G)(A)\) when this identification feeds the comparison theorem.
\end{proof}

\begin{lemma}[Cohomology of the applied resolution as a cosyzygy cokernel]
  \label{mk-lem:cohomology_of_applied_resolution}
  \lean{CategoryTheory.Functor.cohomologyAppliedResolutionIso, CategoryTheory.Functor.gHomologyZeroIso}
  \lean{}
  \dcref{05TA,custom}

  \uses{mk-lem:cosyzygy_ses, mk-lem:applied_cosyzygy_cycles, mk-lem:right_derived_zero_iso_self}
  \textit{Source: Stacks Project, Derived Categories, Tag 05TA
  (\texttt{proposition-enough-acyclics}) and Tag 015D
  (\texttt{lemma-F-acyclic-ses}).}
  With \(G\) additive left exact and the acyclic resolution and cosyzygies of
  Lemma~\ref{mk-lem:cosyzygy_ses}, the cohomology of the applied complex
  \(G(J^\bullet) = \bigl(G(J^0) \to G(J^1) \to \cdots\bigr)\) is computed by
  \[
    H^0\bigl(G(J^\bullet)\bigr) \;\cong\; G(A),
    \qquad
    H^n\bigl(G(J^\bullet)\bigr) \;\cong\;
      \operatorname{coker}\bigl(G(J^{n-1}) \to G(Z^n)\bigr) \quad (n \geq 1),
  \]
  where \(G(J^{n-1}) \to G(Z^n)\) is \(G\) applied to the surjection
  \(J^{n-1} \twoheadrightarrow Z^n\) of \((S_{n-1})\).
\end{lemma}

\begin{proof}
  \uses{mk-lem:cosyzygy_ses, mk-lem:applied_cosyzygy_cycles, mk-lem:right_derived_zero_iso_self}
  The degree-zero statement is the \(n = 0\) case of
  Lemma~\ref{mk-lem:applied_cosyzygy_cycles}: \(H^0(G(J^\bullet)) = \ker(G(J^0) \to G(J^1))
  \cong G(Z^0) = G(A)\). Fix now \(n \geq 1\). By definition of the cohomology of the
  cochain complex \(G(J^\bullet)\) in positive degree,
  \[
    H^n\bigl(G(J^\bullet)\bigr)
      \;=\; \ker\bigl(G(J^n) \to G(J^{n+1})\bigr) \big/
            \operatorname{im}\bigl(G(J^{n-1}) \to G(J^n)\bigr).
  \]
  By Lemma~\ref{mk-lem:applied_cosyzygy_cycles} the numerator is
  \(\ker(G(J^n) \to G(J^{n+1})) \cong G(Z^n)\), realized as the subobject
  \(G(Z^n) \hookrightarrow G(J^n)\). Under this identification the differential
  \(G(J^{n-1}) \to G(J^n)\) lands inside \(G(Z^n)\): it factors as
  \(G(J^{n-1}) \xrightarrow{G(J^{n-1} \twoheadrightarrow Z^n)} G(Z^n)
  \hookrightarrow G(J^n)\) by the factorization of \(J^{n-1} \to J^n\) through its image
  \(Z^n\), so its image equals the image of \(G(J^{n-1}) \to G(Z^n)\) inside the
  subobject \(G(Z^n)\). Therefore the quotient
  \(\ker(G(J^n) \to G(J^{n+1})) / \operatorname{im}(G(J^{n-1}) \to G(J^n))\) is exactly
  \(G(Z^n) / \operatorname{im}(G(J^{n-1}) \to G(Z^n)) =
  \operatorname{coker}(G(J^{n-1}) \to G(Z^n))\). Hence
  \(H^n(G(J^\bullet)) \cong \operatorname{coker}(G(J^{n-1}) \to G(Z^n))\).
\end{proof}

\section{The acyclic-resolution comparison theorem}

\begin{theorem}[An acyclic resolution computes the right-derived functor]\label{mk-lem:acyclic_resolution_computes_derived}
  \lean{CategoryTheory.Functor.rightDerivedIsoOfAcyclicResolution}
  \leanok
  \dcref{015E,custom}

  \uses{mk-def:right_acyclic, mk-lem:cosyzygy_ses, mk-lem:cohomology_of_applied_resolution,
    mk-lem:acyclic_dimension_shift, mk-lem:acyclic_one_iso_coker,
    mk-lem:right_derived_zero_iso_self, mk-lem:right_derived_vanishes_injective}
  \textit{Source: Stacks Project, Derived Categories, Tag 015E (Leray's acyclicity
  lemma) and Tag 05TA (\texttt{proposition-enough-acyclics}).}
  Let \(\mathcal{A}\) be an abelian category with enough injectives, \(\mathcal{B}\)
  abelian, and \(G : \mathcal{A} \to \mathcal{B}\) an additive left-exact functor.
  Let
  \[
    0 \longrightarrow A
      \longrightarrow J^0 \longrightarrow J^1 \longrightarrow J^2
      \longrightarrow \cdots
  \]
  be a resolution of \(A\) (an exact sequence) in which every object \(J^n\) is
  right-\(G\)-acyclic in the sense of Definition~\ref{mk-def:right_acyclic}. Then for
  every \(n \in \mathbb{N}\) the \(n\)-th right-derived functor of \(G\) at \(A\) is
  computed by the resolution: there is an isomorphism in \(\mathcal{B}\)
  \[
    (R^n G)(A) \;\cong\; H^n\bigl(G(J^\bullet)\bigr),
  \]
  where \(G(J^\bullet)\) is the cochain complex
  \(G(J^0) \to G(J^1) \to G(J^2) \to \cdots\) obtained by applying \(G\) degreewise
  to the augmentation-dropped resolution and \(H^n\) is its \(n\)-th cohomology
  object.
\end{theorem}
\begin{proof}
  \leanok
  \uses{mk-def:right_acyclic, mk-lem:cosyzygy_ses, mk-lem:cohomology_of_applied_resolution,
    mk-lem:acyclic_dimension_shift, mk-lem:acyclic_one_iso_coker,
    mk-lem:right_derived_zero_iso_self, mk-lem:right_derived_vanishes_injective}
    The argument is the classical dimension shift, now assembled from the cosyzygy
  infrastructure and the dimension-shift step proved above. Cut the resolution into
  its cosyzygy short exact sequences \((S_m) : 0 \to Z^m \to J^m \to Z^{m+1} \to 0\)
  with \(Z^0 = A\) and acyclic middle terms, as furnished by
  Lemma~\ref{mk-lem:cosyzygy_ses}.

  \emph{Degree zero.} By Lemma~\ref{mk-lem:cohomology_of_applied_resolution},
  \(H^0(G(J^\bullet)) \cong G(A)\), and by \(R^0 G \cong G\)
  (Lemma~\ref{mk-lem:right_derived_zero_iso_self}) the right-hand side is \((R^0 G)(A)\).
  This is the case \(n = 0\).

  \emph{Positive degrees.} Fix \(n \geq 1\). The sequences
  \((S_0), \ldots, (S_{n-2})\) all have acyclic middle terms, so the dimension-shift
  Lemma~\ref{mk-lem:acyclic_dimension_shift} applied to them gives, for each \(k \geq 1\),
  isomorphisms \((R^k G)(Z^{m+1}) \cong (R^{k+1} G)(Z^m)\) for \(0 \le m \le n-2\)
  (with \(Z^0 = A\)). Composing them down the staircase
  \[
    (R^n G)(A)
      \;\cong\; (R^{n-1} G)(Z^1)
      \;\cong\; (R^{n-2} G)(Z^2)
      \;\cong\; \cdots
      \;\cong\; (R^{1} G)(Z^{n-1})
  \]
  reduces \((R^n G)(A)\) to a single first-derived functor. (For \(n = 1\) the
  staircase is empty and \((R^1 G)(A) = (R^1 G)(Z^0)\).) Now apply the lowest-degree
  cokernel description (Lemma~\ref{mk-lem:acyclic_one_iso_coker}) to \((S_{n-1})\),
  \(0 \to Z^{n-1} \to J^{n-1} \to Z^n \to 0\), whose middle term \(J^{n-1}\) is
  acyclic:
  \[
    (R^1 G)(Z^{n-1}) \;\cong\; \operatorname{coker}\bigl(G(J^{n-1}) \to G(Z^n)\bigr).
  \]
  Finally Lemma~\ref{mk-lem:cohomology_of_applied_resolution} identifies this cokernel
  with the cohomology of the applied complex,
  \(\operatorname{coker}(G(J^{n-1}) \to G(Z^n)) \cong H^n(G(J^\bullet))\). Composing
  the staircase, the cokernel description, and this identification gives
  \((R^n G)(A) \cong H^n(G(J^\bullet))\) for every \(n \geq 1\); together with the
  degree-zero case this proves the theorem for every \(n \in \mathbb{N}\).

  Every long exact sequence consumed above is the one produced by the dimension-shift
  Lemma~\ref{mk-lem:acyclic_dimension_shift} and its lowest-degree companion
  Lemma~\ref{mk-lem:acyclic_one_iso_coker}, and is therefore grounded --- through the
  horseshoe lift and the complex-level homology long exact sequence on which those
  lemmas rest --- on the injective-resolution computation of \(R^n G\). The injective
  objects supplied by ``enough injectives'' are themselves right-\(G\)-acyclic
  (Lemma~\ref{mk-lem:right_derived_vanishes_injective}), which is what makes an injective
  resolution a special case of an acyclic resolution and reconciles the present
  formula with the defining computation \((R^n G)(A) \cong H^n(G(I^\bullet))\) for an
  injective resolution \(I^\bullet\).
\end{proof}
